\begin{derivation}[从\eqnrefs{eq:delta-r-definition-dp11}到\eqnrefs{eq:sm-mw-rho-response-dp95}]
把\eqnrefs{eq:delta-r-definition-dp11} 改写为
\begin{equation*}
M_W^2s_W^2=A(1+\Delta r),
\qquad
A\equiv\frac{\pi\alpha}{\sqrt2G_F^{(E)}}.
\end{equation*}
在壳方案中 $s_W^2=1-M_W^2/M_Z^2$，因此定义
\begin{equation*}
f(x)\equiv x\left(1-\frac{x}{M_Z^2}\right),
\qquad x\equiv M_W^2,
\end{equation*}
则精密输入关系就是 $f(x)=A(1+\Delta r)$。树级点 $x_0$ 满足 $f(x_0)=A$。把精确解写成 $x=x(\Delta r)$，满足
\begin{equation*}
 f[x(\Delta r)]=A(1+\Delta r),
 \qquad x(0)=x_0.
\end{equation*}
对 $\Delta r$ 在零点求导，隐函数定理直接给
\begin{equation*}
 f'(x_0)\left.\frac{\dd x}{\dd(\Delta r)}\right|_0=A.
\end{equation*}
而
\begin{equation*}
f'(x_0)
=1-\frac{2M_{W,0}^2}{M_Z^2}
=s_{W,0}^2-c_{W,0}^2,
\qquad
A=x_0s_{W,0}^2.
\end{equation*}
因此平方质量对 $\Delta r$ 的精确线性响应系数为
\begin{equation*}
\left.\frac1{x}\frac{\dd x}{\dd(\Delta r)}\right|_0
=-\frac{s_{W,0}^2}{c_{W,0}^2-s_{W,0}^2}.
\end{equation*}
又因为 $M_W=\sqrt{x}$，链式法则给
\begin{equation*}
\boxed{
\left.\frac1{M_W}\frac{\dd M_W}{\dd(\Delta r)}\right|_0
=-\frac{s_{W,0}^2}{2(c_{W,0}^2-s_{W,0}^2)}.}
\end{equation*}
这就是\eqnrefs{eq:sm-mw-linear-response-dp95} 所表示的一圈线性响应系数；有限 $\Delta r$ 的非线性位移由精确二次方程 $f(x)=A(1+\Delta r)$ 决定。

接着把 $\Delta\rho$ 在 $\Delta r$ 中的系数从在壳反项显式抽出来。树级带电流振幅含组合 $g^2/M_W^2=g_{\rm em}^2/(s_W^2M_W^2)$，因此参数反项中由弱混合角产生
\begin{equation*}
\frac{\delta(g^2/M_W^2)}{g^2/M_W^2}
\supset -2\frac{\delta s_W}{s_W}.
\end{equation*}
而正文\eqnrefs{eq:sw-counterterm-massrelation-dp11} 给
\begin{equation*}
-2\frac{\delta s_W}{s_W}
=\frac{c_W^2}{s_W^2}
\left(\frac{\delta m_W^2}{m_W^2}-\frac{\delta m_Z^2}{m_Z^2}\right).
\end{equation*}
在零动量的守护破缺部分，质量反项差与重整化自能差具有相反号；按正文\eqnrefs{eq:delta-rho-definition-dp11} 的定义，
\begin{equation*}
\left(\frac{\delta m_W^2}{m_W^2}-\frac{\delta m_Z^2}{m_Z^2}\right)_{\rm custodial}
=-\Delta\rho.
\end{equation*}
所以 $\Delta r$ 中这部分必为
\begin{equation*}
\left.\Delta r\right|_{\Delta\rho}
=-\frac{c_W^2}{s_W^2}\Delta\rho.
\end{equation*}
再把光子真空极化主导的电荷运行记为 $\Delta\alpha$，其余有限自能、顶角、盒图与反项归入 $\Delta r_{\rm rem}$，得到
\begin{equation*}
\Delta r
=\Delta\alpha-\frac{c_W^2}{s_W^2}\Delta\rho+\Delta r_{\rm rem}.
\end{equation*}
因此 $-c_W^2/s_W^2$ 不是经验分解系数，而是由在壳定义 $s_W^2=1-m_W^2/m_Z^2$ 对带电流振幅的参数依赖直接固定。
把 $\Delta r$ 中与 $\Delta\rho$ 成线性的部分代入上面的响应系数，得到
\begin{align*}
\left.\frac{\delta M_W}{M_W}\right|_{\Delta\rho}
&=-\frac{s_W^2}{2(c_W^2-s_W^2)}
\left(-\frac{c_W^2}{s_W^2}\Delta\rho\right)\\
&=\frac{c_W^2}{2(c_W^2-s_W^2)}\Delta\rho,
\end{align*}
得到\eqnrefs{eq:sm-mw-rho-response-dp95}。再代入\eqnrefs{eq:delta-rho-top-dp11} 就得到正文的重顶响应\eqnrefs{eq:sm-mw-top-response-dp95}。
\end{derivation}
